\subsection{The second criterion: a rank at one point}

The other way from a dense algebraic locus to an open one is to make the
cycle move, and the classical tool for that is semiregularity. The theorem
below adds nothing to the theory of semiregularity itself. Its content is
that, after \Cref{thm:hdgcount}, \Cref{thm:everyfamily} and
\Cref{thm:closedlocus}, the semiregularity of a single object at a single
point implies the Hodge conjecture for the whole family, with no further
propagation argument. We begin with the object: it exists, and its Chern
character is concentrated in one degree and is a flat Hodge class on the
whole family.

\begin{proposition}[The $K$-theory class at a base point]\label{prop:ktheory}
Let $s_{0}\in\Sigma$ be a base point of \Cref{thm:everyfamily} and write
$A_{0}=A_{s_{0}}$. There are an integer $N\ge1$ and a perfect complex $E$ on
$A_{0}$ with
\[
  \ch_{k}(E)=0\ \ (k\ne n),\qquad \ch_{n}(E)=N\,\omega_{1} .
\]
The tuple $v=\ch(E)$ extends to a flat section of
$\bigoplus_{k}R^{2k}f_{*}\QQ$ over $\cS$ which is of type $(k,k)$ at every
point.
\end{proposition}

\begin{proof}
By \Cref{thm:everyfamily} there is $\xi\in\CH^{n}(A_{0})_{\QQ}$ with
$\cl(\xi)=\omega_{1}$. For a smooth quasi-projective variety over a field the
Chern character $\ch\colon K_{0}(A_{0})_{\QQ}\to\CH^{*}(A_{0})_{\QQ}$ is a
ring isomorphism \cite[Ex.~15.2.16(b) and Cor.~18.3.2]{Ful98}, so there is a
unique $u\in K_{0}(A_{0})_{\QQ}$ with $\ch(u)=\xi$, and its Chern character
has no component outside degree $n$. Choose $N\ge1$ such that $Nu$ is the
image of a class in $K_{0}(A_{0})$; that class is a difference of classes of
vector bundles, hence the class of a perfect complex $E$, and
$\ch(E)=N\xi$. Passing to cohomology gives the displayed identities.

For the last sentence, $\ch(E)$ has a single nonzero component, equal to
$N\omega_{1}$, and $\omega$ is a flat section of type $(n,n)$ everywhere by
\Cref{lem:omegadescends}.
\end{proof}

\begin{theorem}[Buchweitz and Flenner \cite{BF03}; Pridham \cite{Pri24}]\label{thm:BF}
Let $X$ be a smooth projective complex variety, $E$ a perfect complex on $X$,
and $X\hookrightarrow\cX\to\Spec R$ a deformation of $X$ over an Artin local
$\CC$-algebra along which $\ch(E)$ stays of Hodge type. If the semiregularity
map
\[
  \sigma_{E}\colon\Ext^{2}_{X}(E,E)
  \longrightarrow\bigoplus_{q\ge0}H^{q+2}(X,\Omega^{q}_{X}),
  \qquad
  \sigma_{E}(\zeta)=\sum_{q}\tfrac{1}{q!}\,
    \mathrm{tr}\bigl(\zeta\circ\mathrm{At}(E)^{q}\bigr),
\]
is injective, then $E$ extends to a perfect complex on $\cX$.
\end{theorem}

In this generality, for perfect complexes and for every obstruction
annihilated by $\sigma_{E}$ rather than only the curvilinear ones, the
theorem is proved by Pridham \cite{Pri24}, who identifies $\sigma_{E}$ with the
tangent map of a generalised Abel--Jacobi map on the derived moduli of perfect
complexes; for coherent sheaves see also \cite{BLM23}. Perry \cite{Per26}
extends it to noncommutative varieties with group actions, which covers
twisted derived categories and answers a question of Markman.

\begin{theorem}[Semiregularity at one point closes the family]
\label{thm:semiregclosure}
Let $n\ge2$, let $s_{0}$ and $E$ be as in \Cref{prop:ktheory}, and suppose the
semiregularity map $\sigma_{E}$ of \Cref{thm:BF} is injective. Then
$\Sigma=\cD_{n,n}$, and the Hodge conjecture holds in every degree for every
abelian $2n$-fold of $(K,-1,n)$-Weil type with $\Hg(A)=\SU(V,H)$ in that
family.
\end{theorem}

\begin{proof}
Let $\mathfrak{M}$ be the stack of perfect complexes with vanishing negative
self-extensions on the fibres of $f\colon\cA\to\cS$. It is algebraic and
locally of finite presentation over $\cS$ \cite{Lie06}, and at a point $[E]$
over $s_{0}$ its deformation theory relative to $\cS$ has tangent space
$\Ext^{1}(E,E)$ and obstruction space $\Ext^{2}(E,E)$. If
$\Ext^{<0}(E,E)\ne0$, which that stack excludes,\footnote{Lieblich's
stack parametrises the complexes he calls universally gluable, those with
$\Ext^{<0}(E,E)=0$; without that condition a complex can have automorphisms
of negative degree, and the natural moduli object is the derived stack of
To\"en and Vaqui\'e.} use instead the analytic Kuranishi germ of pairs from
the proof of \Cref{thm:p2false}, which has the same tangent and obstruction
spaces; the argument below then gives a nonempty analytic open set in place of
a Zariski open one, which is all that \Cref{cor:allornothing} needs.

By \Cref{prop:ktheory} the Chern character of $E$ stays of Hodge type along
every direction in $\cS$, because it is a flat section of Hodge classes on the
whole base. So \Cref{thm:BF} applies to every infinitesimal deformation of
$A_{0}$ inside the family, and injectivity of $\sigma_{E}$ kills every
obstruction. The morphism $\mathfrak{M}\to\cS$ is therefore formally smooth at
$[E]$, and being locally of finite presentation it is smooth on an open
neighbourhood of $[E]$, which we take connected. Smooth morphisms are open,
so the image of that neighbourhood is a nonempty Zariski open subset
$U\subseteq\cS$.

For $s\in U$ pick a point of the neighbourhood over $s$, that is, a perfect
complex $E_{s}$ on $A_{s}$. Chern characters are locally constant in
families, so $\ch(E_{s})$ is the parallel transport of $\ch(E)$: it is
$N\omega_{1}$ in degree $n$ and zero elsewhere. Chern characters of perfect
complexes are algebraic, so $N\omega_{1}$, and hence $\omega_{1}$, is
algebraic on $A_{s}$; the rest of the Weil line follows by
\Cref{lem:onecomponent}. So $p^{-1}(U)\subseteq\Sigma$, which therefore
contains a nonempty open set, and \Cref{cor:allornothing} gives
$\Sigma=\cD_{n,n}$. Apply \Cref{thm:residual}(a).
\end{proof}

\begin{remark}[The hypothesis is not met by the object above]
\label{rem:semiregunavailable}
Injectivity of $\sigma_{E}$ is automatic for some objects and not for
others, and its degree-zero component already shows both. For a line bundle
$L$ one has $\Ext^{2}(L,L)=H^{2}(\cO_{X})$ and $\sigma_{0}=\mathrm{tr}$ is the
identity on it, so $\sigma_{L}$ is injective for every line bundle on every
variety, and \Cref{thm:BF} returns the classical fact that a divisor class
deforms wherever it stays of type $(1,1)$. The object of \Cref{prop:ktheory}
has rank zero, so its trace component is not an isomorphism, and in general
not injective, the source having no reason to be as small as
$h^{2}(\cO_{A_{0}})=\binom{2n}{2}$. The hypothesis of
\Cref{thm:semiregclosure} therefore rests on the components with $q\ge1$,
which is where semiregularity stops being formal and becomes a computation.

For the object of \Cref{prop:ktheory} the hypothesis fails: by
\Cref{prop:concentrated} it has $\dim\ker\sigma_{E}\ge n^{2}$ at every base
point constructed in this paper (see also \Cref{rem:killsreduction}). The
hypothesis concerns a choice of perfect complex representing a class in
$K_{0}$, and \Cref{prop:concentrated} does not reach every representative.
\Cref{thm:p2forces} and the results after it describe what a representative
satisfying the hypothesis would have to look like, and \Cref{thm:p2false}
shows that there is none. The proof of \Cref{thm:semiregclosure} applies word
for word to the polarised criterion \textup{(P2$'$)} of \Cref{rem:p2prime},
and \Cref{rem:p2primesurvive} records which of the results below carry over
to it.
\end{remark}

\subsection{Consequences of the rank condition}

The hypothesis of \Cref{thm:semiregclosure} is a condition on a rank, which
cannot be inspected until the object is written down. It nevertheless has
consequences that hold for every representative at once, because the
semiregularity map composed with the evaluation map of Hochschild cohomology
is contraction into the Chern character (\Cref{setup:criterion}), and the
Chern character here is a Weil class. The annihilator of a Weil class in
$HT^{2}$ is large, and injectivity forces the evaluation map to vanish on the
whole of it. We compute the annihilator and draw the consequences. They are
the exact analogue for (P2) of what \Cref{thm:candidatefails} and
\Cref{thm:nosplit} do for the secant route in \Cref{sec:secantobjects}, and
they rest on the same local products.

\begin{theorem}[Consequences of injectivity]
\label{thm:p2forces}
Let $A$ be an abelian $2n$-fold of $(K,-1,n)$-Weil type, $n\ge2$, and write
$H^{1}(A,\CC)=V_{+}\oplus V_{-}$ for the eigenspaces of $K$,
$V_{\pm}^{1,0}$, $V_{\pm}^{0,1}$ for their Hodge components and
$T_{\pm}\subset H^{0}(T_{A})=T_{0}A$ for the eigenspaces of $K$ on the tangent
space, each of dimension $n$. Let $\omega=a\alpha_{+}+b\alpha_{-}$ be a
class on the Weil line with $ab\ne0$, for instance any nonzero rational one,
and let $E$ be a nonzero perfect complex on $A$ with $\ch(E)=N\omega$,
$N\ne0$, and no other component.
\begin{enumerate}[label=\textup{(\roman*)},leftmargin=2.6em,itemsep=2pt]
\item The annihilator of $\ch(E)$ in $HT^{2}(A)$ is
  \[
    \Ann_{HT^{2}}(\ch E)
    =\bigl(V_{+}^{0,1}\otimes V_{-}^{0,1}\bigr)
    \;\oplus\;\Hom_{K}\bigl(H^{1,0},H^{0,1}\bigr)
    \;\oplus\;\bigl(T_{+}\otimes T_{-}\bigr),
  \]
  of dimension $n^{2}+2n^{2}+n^{2}=4n^{2}$: the mixed block of $H^{0,2}$,
  which is the annihilator of \Cref{lem:annihilator}; the $K$-linear first
  order deformations of $A$, which contain the tangent space of the polarised
  Weil family; and the mixed bivectors.
\item If $\sigma_{E}$ is injective, then $\operatorname{ev}_{E}$ vanishes on
  this space. Explicitly: $z\cdot\id_{E}=0$ for every $z\in V_{+}^{0,1}\otimes
  V_{-}^{0,1}$; $E$ extends to first order along every $K$-linear deformation
  of $A$; and the translation classes satisfy $A_{\theta}A_{\theta'}=0$ in
  $\Ext^{2}(E,E)$ for every $\theta\in T_{+}$ and $\theta'\in T_{-}$.
\item Under the same hypothesis the algebra map $H^{*}(\cO_{A})\to\Ext^{*}(E,E)$,
  $z\mapsto z\cdot\id_{E}$, kills the ideal generated by the mixed block,
  which is $\bigoplus_{a,b\ge1}\bigwedge^{a}V_{+}^{0,1}\otimes\bigwedge^{b}V_{-}^{0,1}$;
  so the image of $H^{0,k}(A)$ has dimension at most $2\binom{n}{k}$ for
  $1\le k\le n$ and is zero for $k>n$, and every endomorphism $f$ of $E$ has
  $\operatorname{tr}(f)=0$ in $H^{0}(\cO_{A})=\CC$.
\item Under the same hypothesis, if at some point $p$ of its support $E$ is
  locally isomorphic to $i_{*}F$ with $F$ locally free of positive rank on a
  smooth closed subvariety $Z\ni p$ of codimension at least two, then
  $T_{+}\subseteq T_{p}Z$ or $T_{-}\subseteq T_{p}Z$.
\end{enumerate}
\end{theorem}

\begin{proof}
(i) Write $H^{*}(A,\CC)=\bigwedge^{*}(H^{1,0}\oplus H^{0,1})$ and
$\alpha_{\pm}$ for a generator of $\bigwedge^{2n}V_{\pm}
=\bigwedge^{n}V_{\pm}^{1,0}\otimes\bigwedge^{n}V_{\pm}^{0,1}$. The three
summands of $HT^{2}$ act as follows: $z\in H^{0,2}=\bigwedge^{2}H^{0,1}$ by
wedge product; $v\in H^{1}(T_{A})=\Hom(H^{1,0},H^{0,1})$ as the derivation
that replaces one $(1,0)$ factor by its image; $\pi\in H^{0}(\bigwedge^{2}T_{A})
=\bigwedge^{2}(H^{1,0})^{\vee}$ by double contraction of $(1,0)$ factors. For
$z$, the computation is \Cref{lem:annihilator}. For $v$, decompose
$\Hom(H^{1,0},H^{0,1})$ into its four blocks by eigenspaces. An element of
the block $\Hom(V_{+}^{1,0},V_{+}^{0,1})$ replaces a factor of
$\bigwedge^{n}V_{+}^{1,0}$ in $\alpha_{+}$ by an element of $V_{+}^{0,1}$,
which is already exhausted in $\alpha_{+}$, so it kills $\alpha_{+}$; it acts
trivially on $\alpha_{-}$,
which has no $V_{+}^{1,0}$ factor. Likewise $\Hom(V_{-}^{1,0},V_{-}^{0,1})$
kills $\omega$. A nonzero element of the block $\Hom(V_{+}^{1,0},V_{-}^{0,1})$
carries $\alpha_{+}$ to a nonzero element of
$\bigwedge^{n-1}V_{+}^{1,0}\otimes V_{-}^{0,1}\otimes\bigwedge^{n}V_{+}^{0,1}$
and kills $\alpha_{-}$, and a nonzero element of $\Hom(V_{-}^{1,0},V_{+}^{0,1})$
carries $\alpha_{-}$ to a nonzero element of a different summand of
$H^{n-1,n+1}$ and kills $\alpha_{+}$. Since $ab\ne0$ and the two summands are
independent, $v$ kills $\omega$ if and only if its two mixed blocks vanish,
that is, if and only if $v$ is $K$-linear. The space of $K$-linear maps has
dimension $2n^{2}$, and it contains the $n^{2}$-dimensional tangent space of
the polarised family, which is its symmetric part. For $\pi$, a mixed
bivector $\theta\wedge\theta'$ with $\theta\in T_{+}=(V_{+}^{1,0})^{\vee}$ and
$\theta'\in T_{-}=(V_{-}^{1,0})^{\vee}$ kills $\alpha_{+}$ because $\theta'$
kills every factor of $\bigwedge^{n}V_{+}^{1,0}$, and kills $\alpha_{-}$
symmetrically; a nonzero bivector in $\bigwedge^{2}T_{+}$ carries $\alpha_{+}$
to a nonzero element of $\bigwedge^{n-2}V_{+}^{1,0}\otimes\bigwedge^{n}V_{+}^{0,1}$
and kills $\alpha_{-}$, and symmetrically for $\bigwedge^{2}T_{-}$, the two
images lying in different summands; so the annihilator in
$H^{0}(\bigwedge^{2}T_{A})$ is $T_{+}\otimes T_{-}$. Since the images of the
three summands of $HT^{2}$ lie in the three different Hodge components
$H^{n,n+2}$, $H^{n-1,n+1}$ and $H^{n-2,n}$ of $H^{*}(A,\CC)$, the annihilator
is the direct sum of the three annihilators computed, of dimensions $n^{2}$,
$2n^{2}$ and $n^{2}$.

(ii) By the commutative square $\sigma_{E}\circ\operatorname{ev}_{E}
=(\,\cdot\,)\lrcorner\ch(E)$ of \Cref{setup:criterion}, $\sigma_{E}$
vanishes on $\operatorname{ev}_{E}(\Ann_{HT^{2}}(\ch E))$; if $\sigma_{E}$
is injective, $\operatorname{ev}_{E}$ vanishes there. The three explicit
statements are \eqref{eq:evgeneral} on the three summands of (i): on
$H^{0,2}$ the evaluation map is $z\mapsto z\cdot\id_{E}$; on $H^{1}(T_{A})$
it is the obstruction map $v\mapsto v\lrcorner\mathrm{At}(E)$, whose
vanishing on $v$ is the extension of $E$ to first order along $v$
\cite{Tod09}; and on a bivector $\theta\wedge\theta'$ it is
$\tfrac12(\theta\wedge\theta')\lrcorner\mathrm{At}(E)^{2}
=A_{\theta}A_{\theta'}$, as in the proof of \Cref{thm:candidatefails}.

(iii) The map $z\mapsto z\cdot\id_{E}$ is a homomorphism of graded algebras
from $H^{*}(\cO_{A})=\bigwedge^{*}H^{0,1}$ to the graded centre of the Yoneda
algebra, so its kernel is an ideal, and by (ii) it contains the mixed block
$V_{+}^{0,1}\otimes V_{-}^{0,1}$. The ideal generated by the mixed block in
$\bigwedge^{*}(V_{+}^{0,1}\oplus V_{-}^{0,1})
=\bigwedge^{*}V_{+}^{0,1}\otimes\bigwedge^{*}V_{-}^{0,1}$ is the sum of the
summands $\bigwedge^{a}V_{+}^{0,1}\otimes\bigwedge^{b}V_{-}^{0,1}$ with
$a,b\ge1$, since each such summand consists of multiples of the mixed block,
while the remaining summands, with $a=0$ or $b=0$, meet the ideal in zero; the
quotient in degree $k\ge1$ is $\bigwedge^{k}V_{+}^{0,1}\oplus\bigwedge^{k}V_{-}^{0,1}$,
of dimension $2\binom{n}{k}$, and zero for $k>n$. In particular the top class
of $H^{2n}(\cO_{A})$, which is mixed for $n\ge1$, acts by zero. Serre duality
on $A$, whose canonical bundle is trivial, is the perfect pairing
$\Ext^{2n}(E,E)\times\Hom(E,E)\to H^{2n}(\cO_{A})\cong\CC$,
$(\xi,f)\mapsto\operatorname{tr}(\xi\circ f)$; for $\xi=w\cdot\id_{E}$
with $w$ the top class this is $w\cdot\operatorname{tr}(f)$, where
$\operatorname{tr}(f)\in H^{0}(\cO_{A})=\CC$, so $w\cdot\id_{E}=0$ forces
$\operatorname{tr}(f)=0$ for every $f$.

(iv) Let $\cO$ be the local ring of $A$ at $p$ and choose coordinates with
$Z=\{x_{1}=\dots=x_{c}=0\}$, $c\ge2$. Locally $E\cong(\cO/J)^{\oplus r}$
with $J=(x_{1},\dots,x_{c})$, a regular sequence, so
$\Ext^{*}_{\cO}(\cO/J,\cO/J)=\bigwedge^{*}_{\cO/J}N$ as a graded algebra,
where $N=\Hom(J/J^{2},\cO/J)$ is the normal module with basis the classes
$\partial_{1},\dots,\partial_{c}$ dual to $x_{1},\dots,x_{c}$; the
$\Ext^{1}$ is computed on the Koszul resolution exactly as in the proof of
\Cref{lem:localproducts}, and the product is the wedge because the Yoneda
algebra of a Koszul complex is the exterior algebra on the dual of the
generators. For $E$ of rank $r$ the local Yoneda algebra is
$\bigwedge^{*}N\otimes M_{r}(\cO/J)$. By \Cref{lem:edge} the germ of
$A_{\theta}$ at $p$ is the jet class along $\theta$; on the Koszul complex the
jet class of a derivation $D$ is the chain map $-D(d)$, and for
$D=\theta=\sum_{i}c^{i}\partial_{i}$ its component in $\Ext^{1}=N\otimes M_{r}(\cO/J)$
is $\sum_{i\le c}(c^{i}\bmod J)\,\partial_{i}\otimes\id$, the class of
$\theta$ modulo $T_{Z}$. Multiplicativity of the germ map then gives the germ
of $A_{\theta}A_{\theta'}$ as
$(\theta\bmod T_{Z})\wedge(\theta'\bmod T_{Z})\otimes\id$ in
$\bigwedge^{2}N\otimes M_{r}(\cO/J)$, whose value at $p$ is
$\bar\theta(p)\wedge\bar\theta'(p)\in\bigwedge^{2}(T_{p}A/T_{p}Z)$. If
$A_{\theta}A_{\theta'}=0$ for all $\theta\in T_{+}$, $\theta'\in T_{-}$,
then $\bar a\wedge\bar b=0$ for all $\bar a$ in the image $I_{+}$ of $T_{+}$
and $\bar b$ in the image $I_{-}$ of $T_{-}$ in $N_{p}=T_{p}A/T_{p}Z$, and
$I_{+}+I_{-}=N_{p}$ since $T_{+}\oplus T_{-}=T_{p}A$. If both $I_{\pm}$ were
nonzero, every nonzero $\bar b\in I_{-}$ would be proportional to every
nonzero $\bar a\in I_{+}$, so $I_{+}=I_{-}$ would be a line and
$N_{p}=I_{+}+I_{-}$ of dimension one, against $c\ge2$. Hence $I_{+}=0$ or
$I_{-}=0$, which is the claim.
\end{proof}

\begin{corollary}[Locally free components are excluded]\label{cor:p2shape}
Let $A$ be as in \Cref{thm:p2forces} and suppose that no proper abelian
subvariety $B\subset A$ has $T_{0}B\supseteq T_{+}$ or $T_{0}B\supseteq T_{-}$;
equivalently, that $A$ has no nonzero $K$-stable abelian subvariety on which
$K$ acts through a single embedding. This holds whenever $A$ is simple, at
every split member $X\times\widehat X$ with $X$ simple of endomorphism
algebra $\QQ$, and at a general point of the quaternionic locus of
\Cref{setup:quat}.
Let $E$ be a perfect complex on $A$ with $\ch(E)=N\omega$ as above and
$\sigma_{E}$ injective. Then no irreducible component $Z_{0}$ of the support
of $E$ of codimension at least two has a dense open subset along which $E$ is
a locally free sheaf of positive rank on $Z_{0}$. In particular $E$ is not a
sheaf, and along every component of its support of codimension at least two
it has, generically, either at least two nonvanishing cohomology sheaves or a
cohomology sheaf that is not locally free on the reduced support.
\end{corollary}

\begin{proof}
Suppose that $Z_{0}$ is such a component, with $U\subseteq Z_{0}^{\mathrm{sm}}$
dense open and $E|_{U}$ a locally free sheaf of positive rank on $U$. By
\Cref{thm:p2forces}(iv), $U=U_{+}\cup U_{-}$ with
$U_{\pm}=\{p\in U:T_{\pm}\subseteq T_{p}Z_{0}\}$, both closed in $U$; as $U$
is irreducible, $U=U_{+}$ say. Then every constant vector field
$\theta\in T_{+}$ is tangent to $Z_{0}$ along $U$, hence along
$Z_{0}^{\mathrm{sm}}$. The set
$S=\{(p,t)\in Z_{0}\times\CC:p+\exp(t\theta)\in Z_{0}\}$ is a closed
analytic subset of the irreducible space $Z_{0}\times\CC$, and it contains
an open neighbourhood of $Z_{0}^{\mathrm{sm}}\times\{0\}$, because the local
flow of the restriction of $\theta$ to the submanifold $Z_{0}^{\mathrm{sm}}$
exists on such a neighbourhood; so $S=Z_{0}\times\CC$, and
$Z_{0}+\exp(t\theta)\subseteq Z_{0}$, hence $=Z_{0}$, for every $t$. Let
$G=\{x\in A:Z_{0}+x=Z_{0}\}$, a closed algebraic subgroup of $A$, and $B$
its identity component, an abelian subvariety. The connected set
$\exp(T_{+})$ contains $0$ and lies in $G$, hence in $B$, so
$T_{+}\subseteq T_{0}B$. By hypothesis $B=A$, so $Z_{0}=A$, against the
codimension. The last sentence of the statement is the negation of the
condition just excluded. If $E$ were a sheaf, its support would have
codimension at least $n$, because the Chern character of a nonzero sheaf
begins in the codimension of its support, with the support cycle, whose
multiplicities are positive; this is excluded by \Cref{thm:p2support}(iii)
below, whose proof uses nothing from the present one beyond the translation
argument above.

For the equivalence in the hypothesis, let $\sqrt{-d}$ act on $A$ through
an order of $K$. If an abelian subvariety $B\ne A$ has
$T_{0}B\supseteq T_{+}$, then $H_{1}(B,\QQ)\cap\sqrt{-d}\,H_{1}(B,\QQ)$ is a
$K$-stable sub-Hodge structure, hence $H_{1}(C,\QQ)$ for a $K$-stable abelian
subvariety $C\subseteq B$, and $T_{0}C=T_{0}B\cap\sqrt{-d}\,T_{0}B$ still
contains $T_{+}$, on which $\sqrt{-d}$ acts by a scalar. The $K$-stable
complement $C'$ of $C$ for the polarisation is nonzero, since
$C\subseteq B\ne A$, and $T_{0}C'$ is a $K$-stable complement of
$T_{0}C\supseteq T_{+}$, so $T_{0}C'\subseteq T_{-}$ and $K$ acts on $C'$
through a single embedding. Conversely, if a nonzero $K$-stable $C'\ne A$ has
$T_{0}C'\subseteq T_{-}$, its $K$-stable complement $B$ is proper with
$T_{0}B\supseteq T_{+}$. For the three cases: a simple $A$ has no nonzero
proper abelian subvariety; for $X\times\widehat X$ with $X$ simple and
$\End^{0}(X)=\QQ$ the sub-Hodge structures of $H_{1}(X)\otimes\QQ^{2}$ are
$H_{1}(X)\otimes W$ with $W\subseteq\QQ^{2}$, and since $K$ acts on
$\QQ^{2}$ through a quadratic field, $W$ is $K$-stable only for $W=0$ or
$\QQ^{2}$; and at a general quaternionic member $A$ is either simple or, when
the quaternion algebra of \Cref{setup:quat} splits, isogenous to $C\times C$
with $C$ a general
principally polarised abelian $n$-fold, whose sub-Hodge structures are again
$H_{1}(C)\otimes W$ with the same conclusion.
\end{proof}

\subsection{The Hochschild action and the two branches}

\Cref{thm:p2forces} computes the annihilator summand by summand and reads off
what the vanishing of the evaluation map on it means for $E$. That
description hides a structure which is easily brought out and gives more:
the annihilator is the mixed part of a splitting of $HH^{1}(A)$ into two
halves, one for each of the two conjugate pieces $\alpha_{+}$ and
$\alpha_{-}$ of the Weil class, and the same splitting describes the
annihilator in every degree. The results of this subsection concern the
Hochschild cohomology of $A$ alone, and their consequences for $E$ follow
formally.

\begin{setupenv}[Hochschild cohomology in these terms]\label{setup:hh}
For an abelian variety $A$ the Hochschild--Kostant--Rosenberg isomorphism
identifies
\[
  HH^{*}(A)\;\cong\;HT^{*}(A)=\bigoplus_{p,q}H^{p}\bigl(\textstyle\bigwedge^{q}T_{A}\bigr)
  =\textstyle\bigwedge^{*}\bigl(H^{0,1}(A)\oplus H^{0}(T_{A})\bigr),
\]
as graded rings, the grading being $p+q$: the Todd class of $A$ is trivial,
so the twist by $\td^{1/2}$ in \cite{Cal05}, \cite{CRV12} is the identity and
the plain identification is already multiplicative. Write
$HH^{1}(A)=H^{0,1}(A)\oplus H^{0}(T_{A})$, of dimension $4n$, so that
$HH^{*}(A)$ is the exterior algebra on it, of dimension $2^{4n}$. It acts on
$H^{*}(A,\CC)=\bigwedge^{*}\bigl(H^{1,0}\oplus H^{0,1}\bigr)$ by wedging with
the classes of $H^{0,1}$ and contracting the classes of $H^{1,0}$ against
$H^{0}(T_{A})$; the two families of operators involve disjoint generators, so
they pairwise anticommute and the action is an action of the exterior algebra.
In degree two this is the operator $x\lrcorner$ of \Cref{setup:criterion}, and
$\operatorname{ev}_{E}\colon HH^{*}(A)\to\Ext^{*}(E,E)$ is a homomorphism of
graded rings into the graded centre, with
$\sigma_{E}\circ\operatorname{ev}_{E}=(\,\cdot\,)\lrcorner\ch(E)$
\cite[Section 6]{BF08}, \cite{Cal05}, \cite{CRV12}.
\end{setupenv}

\begin{theorem}[The annihilator is the mixed part of a splitting]
\label{thm:hhsplitting}
Let $A$ be of $(K,-1,n)$-Weil type with $n\ge2$, and let
$\omega=a\alpha_{+}+b\alpha_{-}$ lie on the Weil line with $ab\ne0$. Put
\[
  P=\Ann_{HH^{1}}(\alpha_{+})=V_{+}^{0,1}\oplus T_{-},
  \qquad
  Q=\Ann_{HH^{1}}(\alpha_{-})=V_{-}^{0,1}\oplus T_{+} .
\]
\begin{enumerate}[label=\textup{(\roman*)},leftmargin=2.6em,itemsep=2pt]
\item $\dim P=\dim Q=2n$ and $HH^{1}(A)=P\oplus Q$; in particular
  $\Ann_{HH^{1}}(\omega)=0$.
\item $\Ann_{HH^{2}}(\omega)=P\wedge Q$, of dimension $4n^{2}$. This is the
  space computed in \Cref{thm:p2forces}(i), written as the mixed part of the
  splitting.
\item Contraction with $\alpha_{-}$ is injective on $\bigwedge^{*}P$, and
  contraction with $\alpha_{+}$ is injective on $\bigwedge^{*}Q$; the two
  images meet in exactly one line, in degree $2n$.
\item For every $k\ge1$,
  \[
    \Ann_{HH^{k}}(\omega)
    =\Bigl(\bigoplus_{\substack{i+j=k\\ i,j\ge1}}
      \textstyle\bigwedge^{i}P\wedge\bigwedge^{j}Q\Bigr)\;\oplus\;\epsilon_{k},
  \]
  where $\epsilon_{k}=0$ for $k\ne2n$ and $\epsilon_{2n}$ is the line of
  \textup{(iii)}. So
  $\dim\Ann_{HH^{k}}(\omega)=\binom{4n}{k}-2\binom{2n}{k}+\delta_{k,2n}$ for
  $k\ge1$, while $\Ann_{HH^{0}}(\omega)=0$, and the annihilator is the ideal
  generated by its degree-two part together with that one line.
\end{enumerate}
\end{theorem}

\begin{proof}
Choose bases $x_{1},\dots,x_{n}$ of $V_{+}^{1,0}$ and $y_{1},\dots,y_{n}$ of
$V_{-}^{1,0}$, and bases $\bbar x_{i}$ of $V_{+}^{0,1}$ and $\bbar y_{j}$ of
$V_{-}^{0,1}$ (the bar is notation only: complex conjugation carries
$V_{+}^{1,0}$ to $V_{-}^{0,1}$), so that
$\alpha_{+}=x_{1}\cdots x_{n}\bbar x_{1}\cdots\bbar x_{n}$ and
$\alpha_{-}=y_{1}\cdots y_{n}\bbar y_{1}\cdots\bbar y_{n}$, and write
$\partial_{x_{i}},\partial_{y_{j}}$ for the dual basis of
$H^{0}(T_{A})=T_{+}\oplus T_{-}$.

(i) Of the $4n$ basis operators, $\bbar x_{i}\wedge$ kills $\alpha_{+}$
because $\bbar x_{i}$ already occurs in it, and $\partial_{y_{j}}\lrcorner$
kills $\alpha_{+}$ because no $y$ occurs in it; so $P\supseteq
V_{+}^{0,1}\oplus T_{-}$. Conversely, the classes
$\bbar y_{j}\wedge\alpha_{+}$ are linearly independent, as are the classes
$\partial_{x_{i}}\lrcorner\alpha_{+}$, and the two families lie in different
summands of $H^{*}$, so no nonzero combination of the operators
$\bbar y_{j}\wedge$ and $\partial_{x_{i}}\lrcorner$ kills $\alpha_{+}$:
$P=V_{+}^{0,1}\oplus T_{-}$, of dimension $2n$. The computation for $Q$ is the
same with the roles of the two eigenspaces exchanged, and
$P\cap Q=0$ with $\dim P+\dim Q=4n$, so $HH^{1}=P\oplus Q$. A class of degree
one that kills $\omega$ kills $\alpha_{+}$ and $\alpha_{-}$ separately, the
two images lying in different summands and $ab\ne0$, so it lies in
$P\cap Q=0$.

(iii) A monomial $\bbar x_{I}\partial_{y_{J}}\in\bigwedge^{*}P$ sends
$\alpha_{-}$ to $\pm\,\bbar x_{I}\,y_{[n]\setminus J}\,\bbar y_{[n]}$, and
these are linearly independent for distinct pairs $(I,J)$, which is the
injectivity; the statement for $Q$ is symmetric. For the intersection, the
image of $\bigwedge^{*}P$ is spanned by monomials with no factor $x$ and all
$n$ factors $\bbar y$, and the image of $\bigwedge^{*}Q$, by the same
computation, by monomials with no factor $y$ and all $n$ factors $\bbar x$.
Both images are spanned by monomials, so their intersection is spanned by the
monomials common to both, and the only such monomial is
$\bbar x_{[n]}\bbar y_{[n]}$. It is the image of
$\pi_{P}=\bbar x_{[n]}\partial_{y_{[n]}}$ and of
$\pi_{Q}=\bbar y_{[n]}\partial_{x_{[n]}}$, both of degree $2n$. So the
intersection is the single line spanned by that class.

(ii) and (iv) Decompose $\bigwedge^{k}HH^{1}=\bigoplus_{i+j=k}\bigwedge^{i}P
\wedge\bigwedge^{j}Q$. A summand with $i,j\ge1$ annihilates $\alpha_{+}$,
because it has a factor in $P$ and the action is an action of the algebra, and
annihilates $\alpha_{-}$ for the same reason with $Q$; so it lies in the
annihilator of $\omega$. The sum of these summands over all $k$ is the ideal
generated by $P\wedge Q$, since each of their monomials has, after
reordering, a factor $p\wedge q$ with $p\in P$ and $q\in Q$. On the remaining
two summands, $\xi_{P}\in\bigwedge^{k}P$ has
$\xi_{P}\lrcorner\omega=b\,\xi_{P}\lrcorner\alpha_{-}$ and
$\xi_{Q}\in\bigwedge^{k}Q$ has $\xi_{Q}\lrcorner\omega=a\,\xi_{Q}\lrcorner\alpha_{+}$,
so a combination $\xi_{P}+\xi_{Q}$ annihilates $\omega$ if and only if
$b\,\xi_{P}\lrcorner\alpha_{-}=-a\,\xi_{Q}\lrcorner\alpha_{+}$. By (iii) both
contractions are injective and their images meet only in degree $2n$, in one
line; so this happens only for $k=2n$, and then only on the line spanned by
$b^{-1}\pi_{P}-a^{-1}\pi_{Q}$, with $\pi_{P}$ and $\pi_{Q}$ scaled so that
$\pi_{P}\lrcorner\alpha_{-}=\pi_{Q}\lrcorner\alpha_{+}$. In degree zero the
annihilator is zero because $\omega\ne0$. Taking $k=2$ and $n\ge2$, so that
$2\ne2n$, gives (ii), the dimension being
$\dim P\cdot\dim Q=4n^{2}$; the general count is
$\binom{4n}{k}-\binom{2n}{k}-\binom{2n}{k}$ plus the exceptional line.
\end{proof}

\begin{corollary}[Consequences for the Hochschild action]
\label{cor:hhfactor}
Let $A$ and $\omega$ be as above and let $E$ be a nonzero perfect complex on
$A$ with $\ch(E)=N\omega$, $N\ne0$.
\begin{enumerate}[label=\textup{(\roman*)},leftmargin=2.6em,itemsep=2pt]
\item Unconditionally, $\dim\Ext^{2}(E,E)\ge2n(2n-1)$.
\item If $\sigma_{E}$ is injective then $\operatorname{ev}_{E}$ vanishes on
  $P\wedge Q$, hence on the ideal it generates, and the ring homomorphism
  $\operatorname{ev}_{E}\colon HH^{*}(A)\to\Ext^{*}(E,E)$ factors through
  \[
    R\;=\;\textstyle\bigwedge^{*}P\;\oplus_{\CC}\;\bigwedge^{*}Q,
    \qquad \dim_{\CC}R=2^{2n+1}-1,
  \]
  the two exterior algebras glued along their scalars, inside $HH^{*}(A)$ of
  dimension $2^{4n}$. At $n=2$ that is $31$ out of $256$.
\item Granting the compatibility
  $\sigma_{E}\circ\operatorname{ev}_{E}=(\,\cdot\,)\lrcorner\ch(E)$ in degree
  $k$ \cite{Cal05}, \cite{CRV12}, the induced map on $R$ is injective away
  from the line of \Cref{thm:hhsplitting}(iii), so
  \[
    \dim\Ext^{k}(E,E)\;\ge\;2\binom{2n}{k}\quad(1\le k<2n),
    \qquad
    \sum_{k\ge1}\dim\Ext^{k}(E,E)\;\ge\;2^{2n+1}-3 .
  \]
\end{enumerate}
\end{corollary}

\begin{proof}
(i) By \Cref{setup:hh}, $\ker\operatorname{ev}_{E}\subseteq\Ann_{HH^{2}}(\ch E)$
in degree two, since $\sigma_{E}(\operatorname{ev}_{E}(\xi))=\xi\lrcorner\ch(E)$
and the right side is nonzero off the annihilator. Hence
$\dim\Ext^{2}(E,E)\ge\dim HH^{2}(A)-\dim\Ann_{HH^{2}}(\omega)
=\binom{4n}{2}-4n^{2}=2n(2n-1)$ by \Cref{thm:hhsplitting}(ii). No hypothesis
on $\sigma_{E}$ is used.

(ii) $\operatorname{ev}_{E}$ vanishes on the annihilator by
\Cref{thm:p2forces}(ii), its kernel is an ideal because it is a ring map, and
the ideal generated by $P\wedge Q$ in $\bigwedge^{*}(P\oplus Q)
=\bigwedge^{*}P\otimes\bigwedge^{*}Q$ is the sum of the summands with
$i,j\ge1$, whose complement is $\bigwedge^{*}P\oplus_{\CC}\bigwedge^{*}Q$ of
dimension $2\cdot2^{2n}-1$.

(iii) With the compatibility in degree $k$, the argument of (i) gives
$\ker\operatorname{ev}_{E}\subseteq\Ann_{HH^{k}}(\omega)$, and by
\Cref{thm:hhsplitting}(iv) that annihilator meets $\bigwedge^{k}P\oplus
\bigwedge^{k}Q$ only in the exceptional line. So the map induced on $R$ has
kernel contained in that one line and its image has dimension at least
$\dim R-1=2^{2n+1}-2$. One of those dimensions sits in degree zero, where the
image is $\CC\cdot\id_{E}$, nonzero because $E$ is; the remaining
$2^{2n+1}-3$ lie in positive degrees, which is the second inequality. For
$1\le k<2n$ the exceptional line has no component in degree $k$, so
$\bigwedge^{k}P\oplus\bigwedge^{k}Q$, of dimension $2\binom{2n}{k}$, maps
injectively, which is the first.
\end{proof}

\begin{remark}[The two branches]\label{rem:twobranches}
The ring $R$ of \Cref{cor:hhfactor}(ii) is the ring of functions on two odd
affine spaces of dimension $2n$ crossing at the origin. In these terms, if $E$
satisfies the criterion, the deformations of $D^{b}(A)$ along which its
Hochschild action survives form two branches, $P$ and $Q$, which meet only at
the trivial deformation and are the annihilators of the two conjugate pieces
$\alpha_{+}$ and $\alpha_{-}$ of the Weil class. The criterion says that $E$
sees the two halves of the class separately and their interaction not at all.

The splitting $HH^{1}=P\oplus Q$ is not the splitting into eigenspaces of
$K$, which are $V_{+}^{0,1}\oplus T_{+}$ and $V_{-}^{0,1}\oplus T_{-}$; $P$
and $Q$ cut across them. It is the splitting into the annihilators of
$\alpha_{+}$ and of $\alpha_{-}$, and it exists because each of the two
pieces of the Weil class is a top form on an eigenspace. Nor is the
conclusion a contradiction in itself. $R$ is a perfectly good ring, and the
bound of (iii) is consistent with the value of $\chi(E,E)$ computed in
\Cref{prop:objectsize}, since $\sum_{k}(-1)^{k}2\binom{2n}{k}=0$. The
corollary says how large such an $E$ must be and how little of the Hochschild
cohomology it may see, but nothing in the Hochschild computation forbids an
$E$ of the required shape. The exclusion comes from a different direction,
deformation to non-algebraic Weil tori, in \Cref{thm:p2false}.
\end{remark}

\subsection{The numerical form of the criterion and the support}

The same square gives two further results. The first converts the rank of
$\sigma_{E}$ into a single dimension, the form in which a candidate can be
tested. The second locates the support of $E$, and it excludes the first
place one would look: a cycle representing the Weil class.

\begin{theorem}[The criterion as a dimension]\label{thm:p2numerical}
Let $A$, $\omega$ and $E$ be as in \Cref{cor:hhfactor}, so that
$\ch(E)=N\omega$ with $N\ne0$ and no other component.
\begin{enumerate}[label=\textup{(\roman*)},leftmargin=2.6em,itemsep=2pt]
\item If $\dim\Ext^{2}(E,E)=2n(2n-1)$, the least value that
  \Cref{cor:hhfactor}(i) allows, then $\operatorname{ev}_{E}\colon
  HH^{2}(A)\to\Ext^{2}(E,E)$ is surjective with kernel exactly $P\wedge Q$,
  the map $\sigma_{E}$ is injective, and every class of $\Ext^{2}(E,E)$ is a
  sum of products $\operatorname{ev}_{E}(p)\operatorname{ev}_{E}(p')
  +\operatorname{ev}_{E}(q)\operatorname{ev}_{E}(q')$ with $p,p'\in P$ and
  $q,q'\in Q$, hence lies in the graded centre of $\Ext^{*}(E,E)$.
\item Conversely, if $\sigma_{E}$ is injective, then
  $\operatorname{ev}_{E}(HH^{2}(A))$ has dimension exactly $2n(2n-1)$, and
  $\sigma_{E}$ carries it isomorphically onto $HH^{2}(A)\lrcorner\,\omega$.
\item Consequently, if at one base point $s_{0}$ of a family some perfect
  complex $E$ with $\ch(E)=N\omega_{1}$ has $\dim\Ext^{2}(E,E)=2n(2n-1)$,
  then $\Sigma=\cD_{n,n}$, and the Hodge conjecture holds in every degree for
  every member of that family with $\Hg(A)=\SU(V,H)$.
\end{enumerate}
\end{theorem}

\begin{proof}
(i) The square of \Cref{setup:hh} gives $\ker\operatorname{ev}_{E}
\subseteq\Ann_{HH^{2}}(\omega)=P\wedge Q$ in degree two, the equality being
\Cref{thm:hhsplitting}(ii). Hence $\operatorname{rank}\operatorname{ev}_{E}
\ge\binom{4n}{2}-4n^{2}=2n(2n-1)=\dim\Ext^{2}(E,E)$, so
$\operatorname{ev}_{E}$ is onto and its kernel, of dimension $4n^{2}$ and
contained in $P\wedge Q$, is all of $P\wedge Q$. Every class of
$\Ext^{2}(E,E)$ is then $\operatorname{ev}_{E}(\xi)$ for some $\xi$, and
$\sigma_{E}(\operatorname{ev}_{E}(\xi))=N\,\xi\lrcorner\,\omega$ vanishes only
for $\xi\in P\wedge Q=\ker\operatorname{ev}_{E}$; so $\sigma_{E}$ is injective.
Since $HH^{2}(A)=\bigwedge^{2}P\oplus(P\wedge Q)\oplus\bigwedge^{2}Q$ and
$\operatorname{ev}_{E}$ is a ring map into the graded centre, the image is
$\operatorname{ev}_{E}(\bigwedge^{2}P)+\operatorname{ev}_{E}(\bigwedge^{2}Q)$,
which is the description of the classes.

(ii) By \Cref{thm:p2forces}(ii) the kernel of $\operatorname{ev}_{E}$ in
degree two contains $P\wedge Q$, and by the argument of (i) it is contained in
it, so it equals $P\wedge Q$ and the image has dimension $2n(2n-1)$. The
composite $\sigma_{E}\circ\operatorname{ev}_{E}$ is contraction into
$\ch(E)=N\omega$, whose kernel is exactly $P\wedge Q=\ker\operatorname{ev}_{E}$,
so $\sigma_{E}$ maps the image isomorphically onto
$HH^{2}(A)\lrcorner\,\omega$.

(iii) By (i) the map $\sigma_{E}$ is injective, and \Cref{thm:semiregclosure}
applies.
\end{proof}

\Cref{thm:p2numerical}(iii) is a sufficient condition for
$\mathbf{W}(K,n,\delta)$ that mentions no map at all: one object, at one
point, and one number. At $n=2$ the number is $12$ and at $n=3$ it is $30$.
By \Cref{thm:p2false} no complex attains it, but the first two parts of the
theorem hold for the polarised criterion (P2$'$) with the number $r(\gamma)$
of \Cref{thm:p2primenumber}(iv) in place of $2n(2n-1)$, and that is the form
in which the polarised criterion can be tested.

The dimension of $\Ext^{2}(E,E)$ is computable in a way the rank of
$\sigma_{E}$ is not. For a complex built from sheaves whose cohomology is
known, the local-to-global spectral sequence
$E_{2}^{p,q}=H^{p}(A,\cE xt^{q}(E,E))\Rightarrow\Ext^{p+q}(E,E)$ reduces the
count to the cohomology of the local $\cE xt$ sheaves, and on an abelian
variety the cohomology of a sheaf supported on an abelian subvariety, or of a
homogeneous bundle, is a question of linear algebra. The rank of
$\sigma_{E}$, by contrast, requires the trace of the exponential of the
Atiyah class on every class of $\Ext^{2}$, a computation with the object and
not only with its invariants. Part (ii) says that, on the central part of
$\Ext^{2}$, nothing is lost by reading the criterion this way. Whether an
injective $\sigma_{E}$ forces the whole of $\Ext^{2}(E,E)$ to be central is,
by Serre duality, a statement about $\Ext^{2n-2}$, which the square in degree
two does not reach; nothing below uses it.

The second result rests on a local fact about complexes supported at a point,
which we state and prove first.

\begin{lemma}[Two jet classes on a complex of finite length]
\label{lem:finiteproduct}
Let $R=\CC[[x_{1},\dots,x_{c}]]$ with $c\ge2$, let $M$ be a bounded complex of
finitely generated free $R$-modules whose cohomology has finite length, and
put $\chi(M)=\sum_{i}(-1)^{i}\ell(H^{i}(M))$. For $a\in\CC^{c}$ let
$J(a)\in\Ext^{1}_{R}(M,M)$ be the jet class of $M$ along the constant vector
field $\sum_{i}a_{i}\partial_{i}$, represented by the chain map
$-\sum_{i}a_{i}\partial_{i}(d)$ as in the proof of \Cref{thm:p2forces}(iv).
If $\chi(M)\ne0$ and $a,b$ are linearly independent, then $J(a)J(b)\ne0$ in
$\Ext^{2}_{R}(M,M)$.
\end{lemma}

\begin{proof}
Let $T$ be an abelian variety of dimension $c$ and $t\in T$ a point. Linear
coordinates on the universal cover identify the completed local ring
$\widehat{\cO}_{T,t}$ with $R$ in such a way that the translation invariant
vector fields become the constant fields $\sum_{i}a_{i}\partial_{i}$: the
exponential map from $T_{t}T=\CC^{c}$ is a local biholomorphism carrying the
constant vector fields to the translation invariant ones, and its germ at the
origin identifies $\widehat{\cO}_{T,t}$ with $\CC[[x_{1},\dots,x_{c}]]$
compatibly with derivations. A complex of $R$-modules with finite length
cohomology is the same thing as a complex of $\cO_{T,t}$-modules with finite
length cohomology, since modules of finite length over either ring are
modules over $\cO_{T,t}/\mathfrak{m}^{N}$ for $N$ large, so $M$ defines a
perfect complex $M_{T}$ on $T$ supported at $t$. Its local $\cE xt$ sheaves
are skyscrapers at $t$, so the local-to-global spectral sequence gives
$\Ext^{*}_{T}(M_{T},M_{T})=\Ext^{*}_{R}(M,M)$, and by \Cref{lem:edge} the
translation classes $A_{a}$, $A_{b}$ of $M_{T}$ restrict to $J(a)$, $J(b)$.
Apply the square of \Cref{setup:criterion} on $T$ to the bivector $a\wedge b$:
\[
  \sigma_{M_{T}}\bigl(A_{a}A_{b}\bigr)=(a\wedge b)\lrcorner\,\ch(M_{T}).
\]
A sheaf of length $\ell$ supported at a point has Chern character
$\ell\cdot[\mathrm{pt}]$, and Chern characters are additive over the
cohomology of a complex, so $\ch(M_{T})=\chi(M)\,[\mathrm{pt}]$. The class of a
point is a nonzero multiple of $dz_{1}\cdots dz_{c}\,d\bbar z_{1}\cdots
d\bbar z_{c}$, and contracting it with $a\wedge b$ removes two holomorphic
factors and leaves a nonzero class of type $(c-2,c)$, because $a$ and $b$ are
independent. So the right side is nonzero when $\chi(M)\ne0$, hence so is
$A_{a}A_{b}$, and so is its restriction $J(a)J(b)$, the restriction map being
an isomorphism.
\end{proof}

The hypothesis on $\chi(M)$ cannot be dropped. The Euler characteristic is
the only additive invariant here: the Grothendieck group of complexes with
finite length cohomology over a regular local ring is $\ZZ$, generated by the
residue field, since every module of finite length has a composition series
whose factors are all isomorphic to it. The lemma says that the product of two
independent jet classes detects this generator, and on a complex of Euler
characteristic zero the product can vanish. Let $K$ be the Koszul resolution
of the residue field $k$ of $\QQ[x_{1},x_{2}]$, on which the jet class of a
derivation $D$ is the chain map $-D(d_{K})$, and let $u=\partial_{2}(d_{K})
\partial_{1}(d_{K})\colon K\to K[2]$ represent $J(\partial_{2})J(\partial_{1})$.
The entries of $d_{K}$ are $\pm x_{1},\pm x_{2}$, so $u$ has constant
entries. Let $C$ be the cone of $u$, with $C^{m}=K^{m+1}\oplus K^{m+2}$ and
$d_{C}(a,b)=(-d_{K}a,\,u(a)+d_{K}b)$. As $u$ induces zero on cohomology, $C$ has
two cohomology modules, each isomorphic to $k$, so $\chi(C)=0$. Since
$\partial_{i}(u)=0$, the matrix $\partial_{i}(d_{C})$ is block diagonal, and the
product $J(\partial_{2})J(\partial_{1})$ on $C$ is $(a,b)\mapsto(u(a),u(b))$,
which equals $d_{C}H+Hd_{C}$ for $H(a,b)=(b,0)$: it is null homotopic, the
gluing has killed it. On the split complex $k\oplus k[1]$, whose Euler
characteristic is also zero, the same product is $u\oplus u[1]$, which is not
null homotopic because $u\ne0$ in $\Ext^{2}(k,k)$ by the lemma applied to $k$.

\begin{theorem}[The support of an object meeting the criterion]
\label{thm:p2support}
Let $A$ be an abelian $2n$-fold of $(K,-1,n)$-Weil type with $n\ge2$, let
$\omega=a\alpha_{+}+b\alpha_{-}$ with $ab\ne0$, and let $E$ be a perfect
complex on $A$ with $\ch(E)=N\omega$, $N\ne0$, no other component, and
$\sigma_{E}$ injective. For an irreducible component $Z$ of the support of $E$
write
\[
  m_{Z}(E)=\sum_{i}(-1)^{i}\,\ell_{\cO_{A,Z}}\bigl(\cH^{i}(E)_{Z}\bigr)
\]
for the multiplicity of $Z$ in the cycle of $E$.
\begin{enumerate}[label=\textup{(\roman*)},leftmargin=2.6em,itemsep=2pt]
\item If $Z$ has codimension at least two and $m_{Z}(E)\ne0$, then $Z$ is a
  union of cosets of an abelian subvariety $B\subsetneq A$ with
  $T_{0}B\supseteq T_{+}$ or $T_{0}B\supseteq T_{-}$.
\item If $A$ has no proper abelian subvariety $B$ with $T_{0}B\supseteq T_{+}$
  or $T_{0}B\supseteq T_{-}$, which by \Cref{cor:p2shape} holds when $A$ is
  simple, at every split member $X\times\widehat X$ with $X$ simple and
  $\End^{0}(X)=\QQ$, and at a general point of the quaternionic locus, then
  $m_{Z}(E)=0$ for every component $Z$ of codimension at least two, and the
  cycle of $E$ is a divisor.
\item At every member of the family, whatever its endomorphisms, the support
  of $E$ has an irreducible component of codimension less than $n$.
\end{enumerate}
In particular no perfect complex whose support has codimension at least $n$,
at any member of any Weil family, satisfies the criterion of
\Cref{thm:semiregclosure}. This includes every complex assembled from the
structure sheaves of the components of a cycle representing $\omega$, the base
cycle of \Cref{thm:everyfamily} among them.
\end{theorem}

\begin{proof}
(i) Let $c\ge2$ be the codimension of $Z$ and let $U\subseteq Z$ be the dense
open set of smooth points of $Z_{\mathrm{red}}$ that lie on no other component
of the support. Fix $\theta\in T_{+}$ and $\theta'\in T_{-}$; by
\Cref{thm:p2forces}(ii), $A_{\theta}A_{\theta'}=0$ in $\Ext^{2}(E,E)$.
Suppose that at some $p\in U$ the images $\bar\theta,\bar\theta'$ of
$\theta,\theta'$ in $N_{p}=T_{p}A/T_{p}Z$ are linearly independent, and
complete them by images of constant fields $\nu_{3},\dots,\nu_{c}$ to a basis
of $N_{p}$. In linear coordinates on the universal cover the constant fields
are coordinate derivations, and the germ $S$ at $p$ of the translate of the
span of $\theta,\theta',\nu_{3},\dots,\nu_{c}$ is a smooth germ of dimension
$c$, transverse to $Z$ at $p$, to which $\theta$ and $\theta'$ are tangent.
Let $j\colon S\to A$ be the inclusion and $M=Lj^{*}E$ on
$\widehat{\cO}_{S,p}\cong\CC[[x_{1},\dots,x_{c}]]$. Derived restriction is a
ring map $\Ext^{*}(E,E)\to\Ext^{*}(M,M)$, and it carries the germ of
$A_{\theta}$ to the jet class of $M$ along $\theta|_{S}$, because the Atiyah
class is functorial for pullback and $\theta$ is tangent to $S$
\cite[Section 3]{BF03}. So $J(\theta)J(\theta')=0$ on $M$. Near $p$ the support
of $E$ is $Z$, and $Z\cap S=\{p\}$, so $M$ has finite length cohomology; its
Euler characteristic is the intersection multiplicity of the cycle of $E$
with $S$ at $p$, which, $S$ being transverse to $Z$ at a smooth point lying on
no other component, is $m_{Z}(E)$, by Serre's formula for intersection
multiplicities applied to each cohomology sheaf and summed with signs
\cite[Chapter~V]{Ser00}.%
\footnote{Serre's formula expresses $\sum_{i}(-1)^{i}\ell(\mathrm{Tor}_{i}(F,
\cO_{S}))$, for a module $F$ and a smooth germ $S$ meeting its support
properly at $p$, as a sum over the components $Z'$ of the support of
$\ell(F_{Z'})$ times the intersection multiplicity of $Z'$ and $S$ at $p$. For
$S$ transverse to a component smooth at $p$ the intersection multiplicity is
one. For a complex the Tor spectral sequence gives the same formula with
alternating sums over the cohomology, which define $m_{Z}(E)$.}
Since $\theta$ and $\theta'$ are independent directions on $S$,
\Cref{lem:finiteproduct} applied to $M$ gives $m_{Z}(E)=0$.

So if $m_{Z}(E)\ne0$, then at every $p\in U$ and for all $\theta\in T_{+}$,
$\theta'\in T_{-}$ the images satisfy $\bar\theta\wedge\bar\theta'=0$ in
$\bigwedge^{2}N_{p}$. As in the proof of \Cref{thm:p2forces}(iv), the images
$I_{\pm}$ of $T_{\pm}$ in $N_{p}$ span $N_{p}$, which has dimension $c\ge2$,
so one of them is zero: $T_{+}\subseteq T_{p}Z$ or $T_{-}\subseteq T_{p}Z$.
The two conditions cut out closed subsets of the irreducible set $U$ whose
union is $U$, so one of them, say the first, holds on all of $U$, and the
argument of \Cref{cor:p2shape} shows that $Z$ is stable under translation by
the identity component $B$ of its stabiliser, an abelian subvariety with
$T_{0}B\supseteq T_{+}$. Since $Z\ne A$, $B\ne A$.

(ii) follows from (i), since no such $B$ exists; if $A$ itself is a
component of the support, its multiplicity is the rank of $E$, which is zero,
so the cycle of $E$ is a divisor.

(iii) Suppose every component of the support has codimension at least $n$.
Then $\ch_{k}(E)=0$ for $k<n$ and $\ch_{n}(E)=\sum_{Z}m_{Z}(E)[Z]$, the sum
over the components of codimension exactly $n$; this is Riemann--Roch without
denominators in its simplest case \cite[Chapter~15]{Ful98}. By (i), each $Z$
with $m_{Z}(E)\ne0$ is $\pi^{-1}(\pi(Z))$ for the quotient $\pi\colon A\to A/B$
by an abelian subvariety $B$ with $T_{0}B\supseteq T_{+}$ or
$T_{0}B\supseteq T_{-}$, so $[Z]=\pi^{*}[\pi(Z)]$ lies in the subring of
$H^{*}(A,\CC)$ generated by $\pi^{*}H^{1}(A/B,\CC)$, the space of classes of
degree one that vanish on $T_{0}B$ and on its conjugate. When
$T_{0}B\supseteq T_{+}$, a
$(1,0)$-form vanishing on $T_{+}$ lies in $V_{-}^{1,0}$, and a $(0,1)$-form
vanishing on the conjugate of $T_{+}$ lies in $V_{+}^{0,1}$, because complex
conjugation exchanges the two eigenvalues of $K$;%
\footnote{The field $K$ acts on $H_{1}(A,\QQ)$ by rational endomorphisms, so
on the complexification its action commutes with complex conjugation, which
carries the eigenspace of $\sqrt{-d}$ for one embedding of $K$ to that for the
other. On the tangent space the conjugate of $T_{+}$ is the eigenspace in
$T^{0,1}$ for the second embedding; it pairs with $V_{-}^{0,1}$, not with
$V_{+}^{0,1}$.}
so the subring lies in $\bigwedge^{*}(V_{-}^{1,0}\oplus V_{+}^{0,1})$. When
$T_{0}B\supseteq T_{-}$ it lies in $\bigwedge^{*}(V_{+}^{1,0}\oplus
V_{-}^{0,1})$. Choose bases of the four pieces. Both exterior algebras are
spanned by monomials in these bases, so their sum is too, and $\alpha_{+}$, a
single monomial in $n$ generators of $V_{+}^{1,0}$ and $n$ of $V_{+}^{0,1}$,
is not among the monomials spanning it; nor is $\alpha_{-}$. As $ab\ne0$, the
class $N\omega$ has nonzero coordinates on both and is not in the sum, which
contradicts $\ch_{n}(E)=N\omega$. So some component has codimension less
than $n$. The last assertion follows, since a complex assembled from
structure sheaves of the components of a cycle of codimension $n$ is
supported on that cycle.
\end{proof}

\Cref{fig:support} draws the slice and the bivector of the proof.

\begin{figure}[tbp]
\centering
  \includegraphics[width=\textwidth]{fig_support}
\caption{The proof of \Cref{thm:p2support}. The grey sheet stands for a
component of the support of codimension less than $n$, which the theorem
forces to exist; the clay curve $Z$ for a component of codimension at least
two. At a smooth point $p$ of $Z$ the directions $\theta\in T_{+}$ and
$\theta'\in T_{-}$ span the slice $S$, transverse to $Z$, on which $E$ restricts
to a complex of finite length with Euler characteristic $m_{Z}(E)$; the amber
parallelogram is the bivector $\theta\wedge\theta'$, independent modulo
$T_{p}Z$. Carried to a torus of the same dimension, the product of the two
jet classes is seen by the semiregularity map as $m_{Z}(E)$ times the class of
a point contracted with the bivector, which is not zero, so $m_{Z}(E)=0$ or
the bivector degenerates. The teal lines are cosets of an abelian subvariety
$B$ tangent to $T_{+}$, the only components of codimension at least two on
which a nonzero multiplicity survives. The scene is schematic in the
dimensions and exact in the incidences.}
\label{fig:support}
\end{figure}

\subsection{Indecomposability and endomorphisms}

The numerical criterion of \Cref{thm:p2numerical} fixes one number,
$\dim\Ext^{2}(E,E)$. Together with Serre duality, the lower bounds of
\Cref{cor:hhfactor} in the other degrees and the Hodge--Riemann relations for
the Weil class, it fixes a good deal more: the object that carries the class
can be taken indecomposable, and it has at least three linearly independent
endomorphisms at $n=2$ and, granting the compatibility of
\Cref{cor:hhfactor}(iii), at $n=3$, the first case in which the criterion
is needed at a nontrivial discriminant.

\begin{theorem}[The shape of an object meeting the criterion]
\label{thm:p2endo}
Let $A$, $\omega$ and $E$ be as in \Cref{thm:p2numerical}, so that
$\ch(E)=N\omega$ with $N\ne0$ and no other component, and suppose that
$\dim\Ext^{2}(E,E)=2n(2n-1)$. Write $e_{k}=\dim\Ext^{k}(E,E)$.
\begin{enumerate}[label=\textup{(\roman*)},leftmargin=2.6em,itemsep=2pt]
\item $\chi(E,E)=\sum_{k}(-1)^{k}e_{k}=(-1)^{n}N^{2}\int_{A}\omega^{2}$, and
  this number is positive.
\item $E\cong E_{0}\oplus E'$ with $E_{0}$ indecomposable,
  $\ch(E_{0})=\ch(E)$, $\dim\Ext^{2}(E_{0},E_{0})=2n(2n-1)$, $\ch(E')=0$ and
  $\Ext^{2}(E_{0},E')=\Ext^{2}(E',E_{0})=\Ext^{2}(E',E')=0$. So $E_{0}$ meets
  the criterion as well.
\item If $E$ is indecomposable, then $\End(E)=\CC\cdot\id\oplus J$ with $J$
  the nilpotent radical, every $f\in J$ annihilates $\Ext^{2}(E,E)$ and
  $\Ext^{2n-2}(E,E)$ under the Yoneda product, and
  $\operatorname{ev}_{E}(\bigwedge^{2n}P)=\operatorname{ev}_{E}(\bigwedge^{2n}Q)$
  is a single line in $\Ext^{2n}(E,E)$. If moreover $e_{1}=4n$, then $J$
  annihilates $\Ext^{1}(E,E)$ as well.
\item If $n=2$ or $n=3$, then
  \[
    \dim\End(E)\;\ge\;2+\tfrac12\,\chi(E,E)\;>\;2,
  \]
  so $E$ has at least three linearly independent endomorphisms, and at
  $n=3$ also $e_{1}\le e_{0}+9$. In particular no simple object meets the
  criterion: no stable sheaf, no complex with $\End=\CC$, and no image of one
  under an autoequivalence of $D^{b}(A)$; and an indecomposable object that
  meets it has a nonzero nilpotent endomorphism. The smallest profiles
  $(e_{0},\dots,e_{2n})$ that the constraints allow are $(3,8,12,8,3)$ at
  $n=2$, with $\chi=2$, and $(3,12,30,40,30,12,3)$ and $(3,12,30,41,30,12,3)$
  at $n=3$, with $\chi=2$ and $\chi=1$.
\end{enumerate}
Parts (i) and (ii), and part (iv) at $n=2$, use the compatibility
$\sigma_{E}\circ\operatorname{ev}_{E}=(\,\cdot\,)\lrcorner\ch(E)$ of
\Cref{setup:hh} in degree two only, as \Cref{thm:p2numerical} does. Part
(iii), and part (iv) at $n=3$, use it in the further degrees in which
\Cref{cor:hhfactor}(iii) grants it: $2n-2$ and $2n$ for (iii), and three for
(iv).
\end{theorem}

\begin{proof}
(i) The Todd class of $A$ is trivial, so the Riemann--Roch theorem gives
$\chi(E,E)=\int_{A}\ch(E)^{\vee}\ch(E)$, and $\ch(E)^{\vee}=(-1)^{n}N\omega$
since $\omega$ has degree $2n$. The polarisation pairs $V_{+}$ only with
$V_{-}$ (see \Cref{fig:tangent}), so its class $\theta$ lies in
$V_{+}\otimes V_{-}$ and $\theta\wedge\alpha_{\pm}=0$, because $\alpha_{\pm}$
spans the top exterior power of $V_{\pm}$. So $\omega$ is a nonzero rational
primitive class of type $(n,n)$ in the middle degree, and the Hodge--Riemann
relations give $(-1)^{n(2n-1)}\int_{A}\omega^{2}>0$, where
$(-1)^{n(2n-1)}=(-1)^{n}$.

(ii) The category $D^{b}(A)$ has finite dimensional morphism spaces and split
idempotents, so $E=\bigoplus_{i}E_{i}$ with each $E_{i}$ indecomposable. The
evaluation map and the semiregularity map are compatible with direct sums,
and by \Cref{thm:p2numerical}(i) the kernel of $\operatorname{ev}_{E}$ in
degree two is $P\wedge Q$; so $\operatorname{ev}_{E_{i}}(P\wedge Q)=0$ and,
composing with $\sigma_{E_{i}}$, $(P\wedge Q)\lrcorner\ch(E_{i})=0$ for every
$i$. We claim that the only classes of Hodge type killed by $P\wedge Q$ are
those of $\CC\alpha_{+}\oplus\CC\alpha_{-}$. Write
$H^{*}(A,\CC)=\bigwedge^{*}(A_{+}\oplus B_{+}\oplus A_{-}\oplus B_{-})$ with
$A_{\pm}=V_{\pm}^{1,0}$ and $B_{\pm}=V_{\pm}^{0,1}$, and grade it by the four
degrees $(i,j,k,l)$. As in \Cref{setup:hh}, $P$ acts by wedging with $B_{+}$
and contracting $A_{-}$, and $Q$ by wedging with $B_{-}$ and contracting
$A_{+}$, so $P\wedge Q$ is spanned by four kinds of operators, each of which
shifts the four degrees by a fixed amount: wedging with $b_{+}\wedge b_{-}$
kills a component of degree $(i,j,k,l)$ for all choices exactly when $j=n$ or
$l=n$; wedging with $b_{+}$ after contracting $a_{+}$, exactly when $i=0$ or
$j=n$; contracting $a_{-}$ and wedging with $b_{-}$, exactly when $k=0$ or
$l=n$; and contracting $a_{-}$ and $a_{+}$, exactly when $i=0$ or $k=0$. Each
of these is a statement about a tensor product of maps that are injective
unless the degree in question is extreme. A class of type $(p,p)$ has
$i+k=j+l=p$ in every component. If $j=n$, then $l=n$ would force $i=k=n$,
against the fourth condition, so $l<n$ and the third condition gives $k=0$;
hence $i=n+l$, so $i=n$ and $l=0$: the component lies in
$\bigwedge^{n}A_{+}\otimes\bigwedge^{n}B_{+}=\CC\alpha_{+}$. If $j<n$, the
first two conditions give $l=n$ and $i=0$, hence $k=n$ and $j=0$: the
component lies in $\CC\alpha_{-}$. This proves the claim. Since $\ch(E_{i})$ is a rational class of
Hodge type killed by $P\wedge Q$, it lies on the rational Weil plane, and
when it is nonzero \Cref{cor:hhfactor}(i) gives
$\dim\Ext^{2}(E_{i},E_{i})\ge2n(2n-1)$, a nonzero rational Weil class having
both coefficients $a$ and $b$ nonzero. As
$\Ext^{2}(E,E)=\bigoplus_{i,j}\Ext^{2}(E_{i},E_{j})$ has dimension
$2n(2n-1)$ and $\sum_{i}\ch(E_{i})=N\omega\ne0$, exactly one summand $E_{0}$
has nonzero Chern character and every other block of $\Ext^{2}$ vanishes.

(iii) An indecomposable object of a category with finite dimensional
morphism spaces and split idempotents has a local endomorphism ring, so
$\End(E)=\CC\cdot\id\oplus J$ with $J$ nilpotent. By
\Cref{thm:p2numerical}(i), $\operatorname{ev}_{E}$ maps
$\bigwedge^{2}P\oplus\bigwedge^{2}Q$ isomorphically onto $\Ext^{2}(E,E)$.
By Serre duality $e_{2n-2}=e_{2}$, and contraction with $\omega$ is injective
on $\bigwedge^{2n-2}P\oplus\bigwedge^{2n-2}Q$ by \Cref{thm:hhsplitting}(iv),
since $2n-2\ne2n$; so, with the compatibility in
degree $2n-2$, $\operatorname{ev}_{E}$ maps that space isomorphically onto
$\Ext^{2n-2}(E,E)$. In degree $2n$ the semiregularity map is the trace, since
$H^{q+2n}(\Omega^{q}_{A})=0$ for $q>0$, so
$\operatorname{tr}\operatorname{ev}_{E}(\xi)$ is the $H^{0,2n}$ component of
$\xi\lrcorner\ch(E)$; for the generators $\pi_{P}$, $\pi_{Q}$ of
$\bigwedge^{2n}P$ and $\bigwedge^{2n}Q$ it is $Nc_{P}$ and $Nc_{Q}$ with
$c_{P}c_{Q}\ne0$. Indeed $\pi_{P}$ wedges with a basis of $B_{+}$ and contracts
a basis of $A_{-}$, so it kills $\alpha_{+}$, whose factor $\bigwedge^{n}B_{+}$
is already full, and carries $\alpha_{-}$ to $\pm$ the generator
$\bigwedge^{n}B_{+}\wedge\bigwedge^{n}B_{-}$ of $H^{0,2n}$; so $c_{P}=\pm b$
and likewise $c_{Q}=\pm a$, both nonzero. For $f\in\End(E)$ and
$x\in\bigwedge^{2}P\oplus\bigwedge^{2}Q$ write
$f\cdot\operatorname{ev}_{E}(x)=\operatorname{ev}_{E}(\varphi_{f}(x))$; then
$f\mapsto\varphi_{f}$ is a homomorphism of algebras. Let $x\in\bigwedge^{2}P$
and $y\in\bigwedge^{2n-2}P$, and write $x\wedge y=\beta(x,y)\pi_{P}$. Since
$\operatorname{ev}_{E}$ is a ring map that kills the ideal generated by
$P\wedge Q$,
\[
  \beta(x,y)\operatorname{tr}\bigl(f\operatorname{ev}_{E}(\pi_{P})\bigr)
  =\operatorname{tr}\bigl(f\operatorname{ev}_{E}(x)\operatorname{ev}_{E}(y)\bigr)
  =\operatorname{tr}\operatorname{ev}_{E}\bigl(\varphi_{f}(x)\wedge y\bigr)
  =Nc_{P}\,\beta\bigl(\varphi_{f}(x)_{P},y\bigr),
\]
and the pairing $\beta$ is perfect, so $\varphi_{f}(x)_{P}=\mu_{f}x$ with
$\mu_{f}=\operatorname{tr}(f\operatorname{ev}_{E}(\pi_{P}))/Nc_{P}$. The same
computation with $y\in\bigwedge^{2n-2}Q$, where $x\wedge y$ lies in the ideal
generated by $P\wedge Q$, gives $\varphi_{f}(x)_{Q}=0$. So $f$ acts on
$\operatorname{ev}_{E}(\bigwedge^{2}P)$ by a scalar $\mu_{f}$, and in the same
way on $\operatorname{ev}_{E}(\bigwedge^{2}Q)$ by a scalar $\nu_{f}$; $\mu$ and
$\nu$ are characters of $\End(E)$ and vanish on $J$. So $J$ kills
$\Ext^{2}(E,E)$, on both sides because $\operatorname{ev}_{E}$ lands in the
graded centre, and the same argument with the roles of the two degrees
exchanged shows that it kills $\Ext^{2n-2}(E,E)$. In particular
$\operatorname{tr}(f\operatorname{ev}_{E}(\pi_{P}))=\mu_{f}Nc_{P}=0$ for
$f\in J$, so $\operatorname{ev}_{E}(\pi_{P})$ and, likewise,
$\operatorname{ev}_{E}(\pi_{Q})$ lie in the annihilator of $J$ under the
perfect pairing of $\End(E)$ with $\Ext^{2n}(E,E)$. That annihilator is a
line, $J$ having codimension one, and both elements are nonzero, their traces
being $Nc_{P}$ and $Nc_{Q}$; so they span the same line. Finally, if
$e_{1}=4n$ then $\Ext^{1}(E,E)=\operatorname{ev}_{E}(HH^{1})$, because
$\operatorname{ev}_{E}$ is injective on $HH^{1}$, as shown in (iv). For
$f\in J$ and $h\in HH^{1}$ write $f\operatorname{ev}_{E}(h)=
\operatorname{ev}_{E}(h')$; then
$\operatorname{ev}_{E}(h'\wedge h'')=f\operatorname{ev}_{E}(h\wedge h'')=0$ for
every $h''\in HH^{1}$, so $h'\wedge HH^{1}\subseteq P\wedge Q$, which forces
$h'=0$ since $\dim P=\dim Q\ge2$.

(iv) By Serre duality $e_{k}=e_{2n-k}$. The map $\operatorname{ev}_{E}$ is
injective on $HH^{1}$ using degree two alone: if $\operatorname{ev}_{E}(h)=0$
then $\operatorname{ev}_{E}(h\wedge h')=0$ for every $h'$, so
$h\wedge HH^{1}\subseteq\Ann_{HH^{2}}(\omega)=P\wedge Q$, and $h=0$ as in
(iii). So $e_{1}\ge4n$, and \Cref{cor:hhfactor}(iii) gives
$e_{3}\ge2\binom{6}{3}=40$ at $n=3$. At $n=2$ the Euler
characteristic is $2e_{0}-2e_{1}+12$, and at $n=3$ it is
$2e_{0}-2e_{1}+60-e_{3}$. So $2e_{0}=\chi+2e_{1}-12\ge\chi+4$ at $n=2$, and
$2e_{0}=\chi+2e_{1}+e_{3}-60\ge\chi+4$ at $n=3$, where also
$e_{1}=e_{0}+30-\frac12(\chi+e_{3})\le e_{0}+9$. With (i) this gives the
bound. An autoequivalence preserves $\End$, and a stable sheaf is simple.
An indecomposable object with $\dim\End\ge3$ has $J\ne0$. For the smallest
profiles take $e_{0}=3$. At $n=2$, $\chi$ is even and positive, so
$2e_{1}=18-\chi\le16$, and $e_{1}\ge8$ forces $e_{1}=8$ and $\chi=2$. At
$n=3$, $2e_{1}+e_{3}=66-\chi\le65$ with $e_{1}\ge12$ and $e_{3}\ge40$
forces $e_{1}=12$ and $e_{3}\in\{40,41\}$, with $\chi=2$ and $\chi=1$.
\end{proof}

\begin{remark}[Objects excluded by the shape]\label{rem:p2endo}
\Cref{thm:p2endo} gives necessary conditions only, and their effect is to
exclude the objects one would construct first. Line bundles, semi-homogeneous
bundles, skyscraper sheaves, stable sheaves and their Fourier--Mukai
transforms are simple, and direct sums of them are excluded by (ii) together
with \Cref{thm:nosum}. At $n=3$, granting the compatibility of
\Cref{cor:hhfactor}(iii), the shape allowed is an indecomposable object
with a local endomorphism algebra of dimension at least three whose radical
kills $\Ext^{2}$ and $\Ext^{4}$ and pairs to zero with the line
$\operatorname{ev}_{E}(\bigwedge^{6}P)$. The smallest conceivable one has
$\Hom(E,E)$ of dimension three, $\Ext^{1}(E,E)$ equal to the image of
$HH^{1}(A)$, of dimension twelve and killed by the radical, and
$\Ext^{3}(E,E)$ of dimension forty or forty-one. Iterated self-extensions and
thickenings are the natural objects with nilpotent endomorphisms, but an
extension of a simple object $F$ by itself has $\ch=2\ch(F)$, so $F$ would
already have to carry the class, and the question returns one level down.
For $n\ge4$ write $e_{k}=2\binom{2n}{k}+\delta_{k}$ for $1\le k\le2n-1$,
with $\delta_{k}\ge0$ and $\delta_{2}=\delta_{2n-2}=0$. Since
$\sum_{k=1}^{2n-1}(-1)^{k}\binom{2n}{k}=-2$, the Euler characteristic is
$2e_{0}-4+\sum_{k}(-1)^{k}\delta_{k}$, and an excess $\delta_{k}$ in an even
degree $4\le k\le2n-4$ can make it positive with $e_{0}=1$; so the bound on
$\End$ disappears. By
\Cref{thm:p2false} none of these shapes is realised.
\end{remark}
